\begin{proposition}
Пусть $\displaystyle \theta _{n}^{*}( X)$ -- (сильно) состоятельная оценка $\displaystyle \theta $, $\displaystyle \tau ( \theta ) :\mathbb{R}^{k}\rightarrow \mathbb{R}^{s}$ -- непрерывная функция. Тогда $\displaystyle \tau \left( \theta ^{*}\right)$ -- (сильно) состоятельная оценка $\displaystyle \tau ( \theta )$.
\end{proposition}
\begin{proof}
    $\displaystyle \theta _{n}^{*}\xrightarrow{P_{\theta } -a.s.} \theta $. По теореме о наследовании сходимостей $\displaystyle \tau \left( \theta _{n}^{*}\right)\xrightarrow{P_{\theta } -a.s.} \tau ( \theta )$.
    $\displaystyle \theta _{n}^{*}\xrightarrow{P_{\theta }} \theta $. По теореме о наследовании сходимостей $\displaystyle \tau \left( \theta _{n}^{*}\right)\xrightarrow{P_{\theta }} \tau ( \theta )$.
\end{proof}
\begin{lemma}
(о наследовании асимптотической нормальности) Пусть $\displaystyle \theta _{n}^{*}$ -- асимптотически нормальная оценка $\displaystyle \theta \in \Theta \subset \mathbb{R}$ с асимптотической дисперсией $\displaystyle \sigma ^{2}( \theta )$, $\displaystyle \tau ( \theta ) :\mathbb{R}\rightarrow \mathbb{R}$ дифференцируема $\displaystyle \forall \theta \in \Theta $. Тогда $\displaystyle \tau \left( \theta _{n}^{*}( X)\right)$ -- асимптотически нормальная оценка для $\displaystyle \tau ( \theta )$ с асимптотической дисперсией $\displaystyle \sigma ^{2}( \theta ) \cdotp ( \tau '( \theta ))^{2}$.
\end{lemma}
\begin{proof}
    Из определения асимптотической нормальности $\displaystyle \sqrt{n}\left( \theta _{n}^{*}( X) -\theta \right)\xrightarrow{d_{\theta }} \xi \ \sim \ N\left( 0,\ \sigma ^{2}( \theta )\right)$. Обозначим $\displaystyle \xi _{n} =\sqrt{n}\left( \theta _{n}^{*}( X) -\theta \right),\ b_{n} =\frac{1}{\sqrt{n}}$. Тогда по следствию из леммы Слуцкого
    \begin{equation*}
    \frac{\tau ( \theta +\xi _{n} b_{n}) -\tau ( \theta )}{b_{n}} =\sqrt{n}\left( \tau \left( \theta _{n}^{*}\right) -\tau ( \theta )\right)\xrightarrow{d_{\theta }} \xi \cdotp \tau '( \theta ) .
    \end{equation*}
\end{proof}
\begin{exercise}
Показать, что несмещенность оценки не наследуется.
\end{exercise}
\begin{proposition}
(о наследовании асимптотической нормальности в многомерном случае, б/д) Пусть $\displaystyle \Theta \subset \mathbb{R}^{k},\ \tau :\mathbb{R}^{k}\rightarrow \mathbb{R}^{s}$ -- дифференцируема в точках $\displaystyle \Theta $, $\displaystyle \theta _{n}^{*}$ -- асимптотически нормальная оценка для $\displaystyle \theta $ с асимптотической ковариационной матрицей $\displaystyle \Sigma ( \theta )$. Тогда $\displaystyle \tau \left( \theta _{n}^{*}\right)$ -- асимптотически нормальная оценка для $\displaystyle \tau ( \theta )$ с асимптотической ковариационной матрицей $\displaystyle \left(\frac{\partial \tau }{\partial \theta }\right) \cdotp \Sigma ( \theta ) \cdotp \left(\frac{\partial \tau }{\partial \theta }\right)^{T}$.
\end{proposition}